\section{Introduction}

Quantum Federated Learning (QFL) extends the federated learning paradigm by replacing or augmenting classical client models with quantum or hybrid quantum-classical models \cite{fedavg,kairouz2019,qfl_towards_2023,qfl_foundations_2023}. In this work, we study QFL under a privacy-sensitive scenario based on FEMNIST, a benchmark designed for federated settings with heterogeneous and non-i.i.d.\ client data \cite{leaf}. The central objective is twofold: first, to train a global quantum model across five federated clients; second, to investigate machine unlearning when one client must be removed from the federation and its historical contribution should no longer influence the global model.

Machine unlearning is particularly relevant in privacy-preserving learning systems because data subjects or organizations may request removal from a trained model \cite{unlearning_survey_2024}. In a quantum setting, this poses additional challenges due to the structure of variational circuits and the geometry of the parameter space. We therefore adopt Quantum Fisher Information (QFI) as a signal for sensitivity and model influence during the unlearning phase \cite{qfi,federated_unlearning_active_forgetting,forgettable_federated_linear_learning}.

The project is implemented as a clean and reproducible software stack in Python 3.10+, with a central server connected to five clients, a PennyLane-based quantum model, and GPU-aware execution through `lightning.gpu` when available, falling back to `lightning.qubit` or `default.qubit` otherwise \cite{pennylane,lightning2024}. The codebase is organized into separate experiment folders for pure QFL training and QFI-based unlearning, with checkpoints, logs, outputs, and test coverage that support reruns and incremental development.

The experimental design follows two controlled settings. In the first, the model is trained for three federated rounds over five clients. In the second, the federation is trained and then one client, `client_0`, is excluded and the model is retrained over the remaining four clients for three rounds, after which the unlearning signal is summarized using `forget_set_accuracy`, `retain_set_accuracy`, `mia_success_rate`, and `qfi_trace`. The FEMNIST images are reduced to a compact four-feature representation via quadrant compression to make the circuit tractable while preserving the structure of the federated pipeline.

The paper is organized as follows. Section~\ref{sec:background} summarizes the background on federated learning, QFL, PennyLane, and QFI. Section~\ref{sec:methodology} introduces the architecture and algorithmic design. Section~\ref{sec:experiments} details the experimental protocol. Section~\ref{sec:results_expected} defines the expected outputs and evaluation metrics. Section~\ref{sec:discussion} discusses limitations and research directions.
